\subsection{张量空间中的语言模型}
\label{sec:lmts}


% \subsubsection*{\textbf{Text representation}}

% Suppose a given text consists of $N$ words $ X = \{x^{(1)}, x^{(2)}, \cdots ,x^{(N)}\}$ and let $\vf_i$ be a feature extractor,  each view of $X$ could be represented as a concatenation of feature maps:

% $$
% \small
% \Phi(X)_\textrm{concat} = (\vf_1 (x^{(1)});  \vf_2 (x^{(2)}); \cdots ;  \vf_N (x^{(N)}))
% $$

% where $\vf_i$ could be one-hot encoding or word embedding.

自然语言中的特征（例如 token 或词）之间通常存在复杂的依赖关系 \citep{hou2013mining}\footnote{这类依赖关系（包括搭配）曾被看作纠缠的一种类比 \citep{hou2013mining}。}，而诸如特征拼接等标准方法难以很好地刻画这些依赖。在因子分解机（factorization machines）中，任意特征之间也存在类似的交互 \citep{rendle2010factorization}。给定由 $N$ 个词构成的文本 $X = [x^{(1)}, x^{(2)}, \cdots ,x^{(N)}]$ 以及特征提取器 $\vf_i \in \mathbb{R}^{I_i}$（它可以是独热编码（one-hot encoding）或词嵌入），我们现在定义一种旨在刻画这些依赖关系的 $X$ 的表示：
%
\begin{equation}
\begin{split}
    \label{eq: Phi X}
    \Phi(X)
    &= \vf_1 (x^{(1)}) \otimes  \vf_2 (x^{(2)}) \cdots \otimes \vf_N (x^{(N)}) \\
    &= \bigotimes^N_{i=1} \vf_i (x^{(i)}) \\
    % &= \includegraphics[width=2cm, height = 0.9cm]{figure/G.png}
\end{split}
\end{equation}
%
其中张量空间（tensor space）为 $\mathbb{R}^{I_1}\otimes\mathbb{R}^{I_2} \otimes\cdots\otimes\mathbb{R}^{I_N}$。$\vf_i$ 的每个分量表示独立的承载含义的单元，例如词素或隐因子。为简单起见，我们在后续章节中假设文本共享同一个独热编码 $\vf(x^{(t)}) \in \mathbb{R}^{|V|}$。于是，$\Phi(X)$ 是一个 $|V|^N$ 维的张量，它记录了 $X$ 中词的所有可能组合。

受 \cite{zhang2019generalized, kossaifi2020tensor} 启发，我们定义一个张量回归模型来计算每个文本 $X$ 的概率：
%
\begin{align}
    \label{eq: general definition}
    p(X)
    &= \langle\mathcal{A}, \Phi(X)\rangle \nonumber \\
    &= \sum_{i_1,i_2,\cdots,i_N = 1}^{|V|} \mathcal{A}_{i_1,\cdots, i_N} \cdot \Phi(X)_{i_1,\cdots, i_N}
\end{align}
%
其中 $\langle \cdot \rangle$ 表示两个同尺寸张量的内积，而 $\mathcal{A}$ 是一个回归权重张量，它与张量空间 $\mathbb{V}^{\otimes N} = \underbrace{\mathbb{V} \otimes \cdots \otimes \mathbb{V}}_N $（其中 $\mathbb{V}$ 指 $\mathbb{R}^{|V|}$）中的 $\Phi(X)$ 具有相同的形状。类似的函数曾在 \cite{novikov2016exponential, NIPS2016_6211, khrulkov2018expressive, zhang2019generalized} 中被考虑过。

% The diagrammatic notation of Eq. \ref{eq: general definition} is \includegraphics[scale=0.10]{figure/d_original.png}.

